\chapter{Discussion}

\section{Interpretation of Results}

The 4,080-run benchmark establishes three principal findings. First, the simulation engine consistently outperforms the naive baseline across all primary metrics, with the improvement being statistically significant (all $p < 0.001$) and practically meaningful (Cohen's $d$ ranging from $-0.41$ to $+1.14$). Second, the engine achieves universal advantage across all 20 benchmark scenarios, demonstrating that the improvement is not an artifact of scenario selection but a general property of the engine architecture. Third, outcome fidelity exhibits the largest engine--naive gap, suggesting that the structured action layer contributes more to simulation realism than persona drift reduction alone.

The magnitude of the persona drift improvement merits careful interpretation. A 5.7\% reduction in drift (from 0.178 to 0.168 MAE) translates to an average per-trait error reduction of approximately 0.010 on a 0--1 scale. While this may appear small, the OCEAN scale is inherently compressed in LLM-generated text: the dynamic calibration formula (Eq.~\ref{eq:dynamic}) centers trait estimates at 0.50 with calibration coefficients of 0.50--0.65, meaning that the effective range of estimated trait values is approximately $[0.20, 0.80]$. Within this range, a 0.010 reduction represents a more substantial improvement than the raw number suggests.

The outcome fidelity result (87.6\% improvement, Cohen's $d = 1.14$) is the most striking finding. The engine's action layer---which extracts structured actions, tracks world-state deltas, and validates action-plan alignment---transforms simulation dialogue from unconstrained conversation into traceable decision-making. The naive baseline, lacking any action extraction, produces dialogue that may sound reasonable but generates no structured outcomes comparable to real-world deliberation results. This finding suggests that persona fidelity and outcome fidelity, while correlated, are partially independent: the action layer improves outcome realism through mechanisms (structured action types, world-state tracking) that operate independently of personality trait management.


\section{RLHF Sycophancy as a Structural Constraint}

The most important finding of this study is not the engine's improvement over the baseline but rather the identification of a structural ceiling imposed by RLHF training on behavioral fidelity. This ceiling manifests in three quantifiable ways.

\textbf{Neuroticism suppression.} The \texttt{self\_doubt\_count} feature records 0.000 across all 4,080 runs under both conditions. LLMs trained with RLHF refuse to generate expressions of self-doubt, anxiety, or emotional vulnerability, regardless of the assigned Neuroticism prior. The N-trait absolute error ($\approx$0.170) is identical across both conditions, confirming that this is an LLM-level limitation rather than an engine-level failure. This finding is consistent with RLHF's optimization objective: human evaluators penalize outputs that express uncertainty or emotional distress, creating a training signal that suppresses these behaviors. The engine, operating at the selection and orchestration level, cannot produce what the underlying model refuses to generate.

\textbf{Conscientiousness estimation floor.} The C-trait absolute error ($\approx$0.222) is the highest of all five traits under both conditions. While the dynamic calibration technique reduced C-error by 38\% relative to static estimation in earlier experiments, the remaining error appears to represent a measurement floor: in policy deliberation contexts, all actors use planning language, structure markers, and action items, making it difficult to distinguish genuinely high-C actors from actors using context-appropriate language. This is a domain-specific limitation rather than an RLHF constraint---the same actors in a casual social conversation might exhibit more differentiated C expression.

\textbf{De-escalation bias.} The outcome analysis reveals that 100\% of adversarial scenarios are simulated as non-adversarial, with mean tension of 0.000 in expected-high-tension contexts. This represents a wholesale failure to simulate conflict---a core requirement for realistic policy deliberation. The engine's sycophancy detection (Algorithm~\ref{alg:sycophancy}) reduces sycophantic acknowledgments but cannot prevent the underlying model from generating fundamentally cooperative language. An actor assigned an adversarial disposition with high disposition strength still produces hedged, collaborative responses because the RLHF training signal overrides the persona prompt.

These three manifestations of the RLHF ceiling suggest that substantial improvements in persona fidelity for suppressed traits will require intervention at the model level (fine-tuning or Direct Preference Optimization) rather than the orchestration level. The engine has effectively maximized what can be achieved through prompt engineering and candidate selection; further progress requires changing the probability distribution from which candidates are sampled.


\section{Effectiveness of Individual Components}

While the present benchmark compares the full engine against the naive baseline rather than conducting component-level ablation, the iterative development history (Experiments~1--6, Arms~A/B) provides insight into each component's contribution.

\textbf{Commitment Lifecycle.} The elimination of contradictions (0.000 across all 4,080 runs) validates the soft prompt injection approach. The informational design---presenting ``Your prior positions:'' without directive instructions---achieves zero contradictions while maintaining behavioral diversity, a result that the directive variant (``stay consistent unless explicitly revising'') could not achieve due to diversity collapse (0.45). This finding has broader implications for LLM prompt design: declarative state presentation outperforms imperative behavior commands for maintaining consistency without sacrificing flexibility.

\textbf{4-Slot Candidate Pool.} The 16.5\% reduction in Extraversion error and 10.0\% reduction in Openness error demonstrate that candidate pooling is most effective for traits with clear psycholinguistic signals. When the scoring function can meaningfully discriminate among candidates (as with word count and idea generation), multiple candidates provide genuine selection value. For traits with weaker signals (Conscientiousness, Neuroticism), the four candidates tend to score similarly, reducing the pool's marginal value.

\textbf{Dynamic OCEAN Calibration.} The relative-to-conversation-baseline approach (Eq.~\ref{eq:dynamic}) is essential for accurate trait estimation in multi-speaker contexts. This technique reduced C-error by 38\% compared to static coefficient approaches in prior experiments. Its effectiveness stems from the normalization of domain-specific language patterns: absolute feature counts are uninformative when all actors in a policy discussion use planning language, but relative-to-baseline counts correctly identify actors who use \textit{more} planning language than their peers.

\textbf{Sycophancy Detection.} The multi-layered detection system (Algorithm~\ref{alg:sycophancy}) reduces sycophantic acknowledgments but does not eliminate them. Across the benchmark, the engine reduces acknowledgment-based sycophancy relative to the naive baseline, but the underlying model's cooperative bias remains the dominant force. The scenario polarity amplification (Layer~2) and strategic disposition adjustment (Layer~3) provide meaningful sensitivity calibration, but their effect is bounded by the candidate pool's limited diversity on sycophancy-related features.

\textbf{Action Layer.} The 12-action vocabulary with world-state deltas is the primary driver of the outcome fidelity improvement. Action-plan alignment of 0.968 indicates that actors' actions are highly consistent with their stated plans, and planned action coverage of 0.928 indicates that most planned actions are eventually executed. However, structured action validity (0.641) reveals that approximately 36\% of extracted actions fail validation, primarily due to phase-gate restrictions (OPENING operates in shadow mode) and unresolvable action targets.


\section{Iterative Refinement Validation}

A follow-up 240-run benchmark (1,080 total runs including both conditions, described as ``iterative refinement'') was conducted with identical scenarios and conditions but with hyperparameter calibration informed by the initial benchmark results. Specifically, relationship delta clamping and fuzzy action validation were adjusted based on diagnostic findings.

The iterative refinement results confirm the main benchmark findings: the engine achieves drift of 0.1685 ($\pm$0.016) vs.\ naive drift of 0.1776 ($\pm$0.016), with 95\% CIs of $[0.167, 0.170]$ and $[0.176, 0.179]$ respectively. Envelope violations are 2.032 vs.\ 2.255, and commitment fulfillment is 0.726 vs.\ 0.646. Per-trait patterns are identical: O ($-11.2\%$) and E ($-15.1\%$) show improvement, while C ($-0.9\%$) and N (+0.9\%) show no change. Outcome fidelity is 0.290 vs.\ 0.161 (Cohen's $d = 1.097$).

The consistency between the 4,080-run main benchmark and the 1,080-run iterative refinement provides strong evidence that the engine's advantages are robust and reproducible. The near-identical per-trait patterns confirm that the RLHF ceiling is a stable structural feature rather than a sampling artifact.


\section{Limitations}

Several limitations qualify the interpretation of these results.

\textbf{Feature-based evaluation.} The 30-feature behavioral extraction pipeline uses keyword matching and pattern detection rather than human judgment. While this approach ensures reproducibility and eliminates evaluator bias~\cite{stureborg2024inconsistent}, it may miss nuanced behavioral signals that human evaluators would detect. The features are designed as proxies for OCEAN trait expression, not direct measurements of personality. A human evaluation study comparing feature-based and human assessments of persona fidelity would strengthen the validity of these metrics.

\textbf{Single candidate model.} All simulations use DeepSeek~V3 as the text generation model. The RLHF behavioral ceiling observed in this study may differ in magnitude and trait specificity across models: Claude, GPT-4, and Gemini undergo different RLHF procedures with different safety calibrations. A cross-model comparison would clarify which aspects of the fidelity ceiling are universal to RLHF and which are model-specific.

\textbf{Adversarial scenario failure.} The 100\% de-escalation bias represents a fundamental limitation. Policy deliberation frequently involves genuine conflict, strategic deception, and power-asymmetric negotiations. The engine's inability to simulate adversarial dynamics limits its applicability to scenarios where stakeholder cooperation is the expected mode of interaction.

\textbf{Outcome fidelity approximation.} Outcome fidelity is measured through automated comparison against documented real-world outcomes, not expert validation. The six-dimensional scoring rubric (resolution accuracy, stakeholder positioning, dynamics patterns, turning points, world-state changes, action distribution) provides structured evaluation, but the ground truth itself is subject to interpretation---real-world policy outcomes are complex, multi-causal events that resist binary characterization.

\textbf{Conscientiousness estimation.} The C-trait error floor ($\approx$0.222) may partially reflect a measurement limitation rather than a simulation limitation. In policy deliberation contexts where planning language is ubiquitous, the feature-based approach may lack the sensitivity to distinguish between context-driven and trait-driven Conscientiousness expression.


\section{Comparison with Related Work}

The present study's contributions can be situated relative to three categories of prior work.

\textbf{Persona simulation approaches.} PersonaLLM, CharacterLLM, and RoleLLM address persona fidelity in single-turn or two-party dialogue settings. The present study extends this to multi-stakeholder group deliberation, where the challenge is qualitatively different: agents must maintain individual personas while responding to social pressure from multiple interlocutors simultaneously. The 4-Slot Candidate Pool and Sycophancy Detection components are specifically designed for this multi-party challenge and have no analogue in single-agent persona systems.

\textbf{Multi-agent simulation frameworks.} Generative Agents~\cite{park2023generative} and AgentSociety demonstrate emergent social behaviors at scale but do not measure individual behavioral fidelity. The present study's 30-feature extraction pipeline and OCEAN estimation framework provide the first quantitative methodology for measuring turn-by-turn persona fidelity in multi-agent interaction. The engine's paired experimental design (same script, both conditions) establishes a rigorous evaluation framework that could be adopted by future multi-agent simulation studies.

\textbf{Policy simulation benchmarks.} PolicySimEval establishes that current systems achieve only 24.5\% accuracy on policy simulation tasks. While the present study's outcome fidelity metric (0.302 for the engine) is not directly comparable due to different scoring rubrics, both results confirm that policy simulation remains a substantially unsolved problem. The present study's contribution is to demonstrate that engine-level persona fidelity control provides a statistically significant improvement over uncontrolled generation, establishing a reproducible baseline for future work.
